% Part: first-order-logic
% Chapter: models-of-arithmetic
% Section: models-of-pa

\documentclass[../../../include/open-logic-section]{subfiles}

\begin{document}

\olfileid[hi]{mod}{mar}{mpa}
\section{$\Th{PA}$ के प्रतिरूप}

\begin{explain}
~$\Th{TA}$ का प्रत्येक अमानक प्रतिरूप~$\Th{PA}$ का भी अमानक प्रतिरूप है।
हम जानते हैं कि~$\Th{TA}$ के, और इसलिए~$\Th{PA}$ के, अमानक प्रतिरूप
अस्तित्व में हैं। हम यह भी जानते हैं कि ऐसे अमानक प्रतिरूपों में अमानक
“संख्याएँ” होती हैं, अर्थात् प्रांत के ऐसे !!{element}s जो सभी मानक
“संख्याओं” के “आगे” होते हैं। किंतु वे किस प्रकार व्यवस्थित हैं? वे कितने
हैं? हमने देखा है कि दुर्बलतर सिद्धांत~$\Th{Q}$ के प्रतिरूपों में केवल एक
अमानक संख्या भी हो सकती है। किंतु ये सरल !!{structure}s न~$\Th{PA}$ की
प्रतिरूप हैं, न~$\Th{TA}$ की।

$\Th{PA}$ या $\Th{TA}$ के प्रतिरूपों की संरचना समझने की कुंजी यह देखना है
कि इन सिद्धांतों में कौन-से तथ्य !!{derivable} हैं। उदाहरणार्थ, $\Th{PA}$
पहले ही सिद्ध करता है कि $\lforall[x][\eq/[x][x']]$ और
$\lforall[x][\lforall[y][\eq[(x+y)][(y+x)]]]$; अतः यह सरल संरचनाओं को
(जिनमें ये !!{sentence}s असत्य हैं)~$\Th{PA}$ के प्रतिरूप होने से रोकता है।

मान लें कि~$\Struct{M}$,~$\Th{PA}$ का प्रतिरूप है। तब यदि
$\Th{PA} \Proves
!A$, तो $\Sat{M}{!A}$। फिर से $\nszero$ का उपयोग
$\Assign{\Obj{0}}{M}$ के लिए, $\nssucc$ का $\Assign{\prime}{M}$ के लिए,
$\nsplus$ का $\Assign{+}{M}$ के लिए, $\nstimes$ का
$\Assign{\times}{M}$ के लिए, और $\nsless$ का $\Assign{<}{M}$ के लिए करें।
तब कोई भी !!{sentence}~$!A$, $\nszero$,
$\nssucc$, $\nsplus$, $\nstimes$ और $\nsless$ के बारे में कोई शर्त
बताता है; यदि $\Sat{M}{!A}$, तो उस शर्त को संतुष्ट होना चाहिए। उदाहरणार्थ,
यदि $\Sat{M}{!Q_1}$, अर्थात्
$\Sat{M}{\lforall[x][\lforall[y][(\eq[x'][y'] \lif \eq[x][y])]]}$,
तो $\nssucc$ को !!{injective} होना चाहिए।
\end{explain}

\begin{prop}
~$\Struct{M}$ में $\nsless$ एक रैखिक सख्त क्रम है; अर्थात् वह निम्न को
संतुष्ट करता है:
\begin{enumerate}
\item $x \nsless x$ नहीं, किसी भी~$x \in \Domain{M}$ के लिए।
\item यदि $x \nsless y$ और $y \nsless z$, तो $x \nsless z$।
\item किसी भी $x \neq y$ के लिए, $x \nsless y$ या $y \nsless x$।
\end{enumerate}
\end{prop}

\begin{proof}
$\Th{PA}$ निम्न को सिद्ध करता है:
\begin{enumerate}
\item $\lforall[x][\lnot x < x]$
\item $\lforall[x][\lforall[y][\lforall[z][((x < y \land y < z) \lif x < z)]]]$
\item $\lforall[x][\lforall[y][((x < y \lor y < x) \lor \eq[x][y]))]]$
\end{enumerate}
\end{proof}

\begin{prop}
\ollabel{prop:M-discrete} $\nszero$,~$\Domain{M}$ का $\nsless$-क्रम में
सबसे छोटा !!{element} है। किसी भी $x$ के लिए $x \nsless
x^\nssucc$, और
$x^\nssucc$ इस गुण वाला $\nsless$-सबसे छोटा !!{element} है। किसी भी $x$
के लिए एक अद्वितीय $y$ ऐसा है कि $y^\nssucc = x$। (हम~$y$ को
~$x$ का,~$\Struct{M}$ में, “पूर्ववर्ती” कहते हैं और इसे~$^\nssucc x$ से
निरूपित करते हैं।)
\end{prop}

\begin{proof}
अभ्यास।
\end{proof}

\begin{prob}
$\Lang{L_A}$ में ऐसे !!{sentence}s खोजिए जो~$\Th{PA}$ में
!!{derivable} हों (और इसलिए~$\Struct{N}$ में सत्य हों) तथा
\olref[mod][mar][mpa]{prop:M-discrete} में $\nszero$, $\nssucc$ और
$\nsless$ के गुणों को सुनिश्चित करें।
\end{prob}

\begin{prop}
~$\Struct{M}$ के सभी मानक !!{element}s, सभी अमानक !!{element}s से
(~$\nsless$ के अनुसार) छोटे हैं।
\end{prop}

\begin{proof}
हम $n$ का उपयोग $\Value{\num{n}}{M}$ के संक्षिप्त रूप में करेंगे, जो
~$\Struct{M}$ का मानक !!{element} है। $\Th{Q}$ पहले ही सिद्ध करता है कि
किसी भी~$n
\in \Nat$ के लिए,
$\lforall[x][(x < \num{n}' \lif (\eq[x][\num{0}] \lor
  \eq[x][\num{1}] \lor \dots \lor \eq[x][\num{n}]))]$। कोई भी
!!{element}, $\nsless \nszero$ नहीं है। इसलिए यदि $n$ मानक और $x$ अमानक
हो, तो $x \nsless n$ नहीं हो सकता। परिभाषा से, अमानक अवयव वह है जो
$\Value{\num{n}}{M}$ नहीं है किसी भी~$n
\in \Nat$ के लिए, इसलिए $x \neq n$ भी।
चूँकि $\nsless$ एक रैखिक क्रम है, $n \nsless x$ होना आवश्यक है।
\end{proof}

\begin{prop}
प्रत्येक अमानक !!{element}~$x$,~$\Domain{M}$ के निम्न उपसमुच्चय का एक अवयव
है:
\[
\dots ^{\nssucc\nssucc\nssucc}x \nsless ^{\nssucc\nssucc}x \nsless
^{\nssucc}x \nsless x \nsless x^{\nssucc} \nsless x^{\nssucc\nssucc}
\nsless x^{\nssucc\nssucc\nssucc} \nsless \dots
\]
हम इस उपसमुच्चय को \emph{$x$ का खंड} कहते हैं और इसे $[x]$ लिखते हैं।
इसमें न कोई सबसे छोटा !!{element} है, न कोई सबसे बड़ा। इसे उन
$y \in \Domain{M}$ के समुच्चय के रूप में अभिलक्षित किया जा सकता है जिनके
लिए किसी मानक~$n$ पर $x \nsplus n = y$ या $y \nsplus n = x$।
\end{prop}

\begin{proof}
स्पष्टतः ऐसा समुच्चय~$[x]$ सदैव अस्तित्व में है, क्योंकि प्रत्येक
!!{element}~$y$,~$\Domain{M}$ में, एक अद्वितीय परवर्ती~$y^\nssucc$ और एक अद्वितीय
पूर्ववर्ती~$^\nssucc y$ है। क्रमागत !!{element}s $y$, $y^\nssucc$ के लिए
$y \nsless y^\nssucc$, और $y^\nssucc$ वह $\nsless$-सबसे छोटा
!!{element} है~$\Domain{M}$ का, जिससे $y$, $\nsless$-छोटा है। चूँकि
सदैव $^\nssucc y \nsless y$ और $y \nsless y^\nssucc$, इसलिए $[x]$ में
न कोई सबसे छोटा !!{element} है, न सबसे बड़ा। यदि $y \in [x]$, तो
$x \in [y]$, क्योंकि तब या तो $y^{\nssucc\dots\nssucc} = x$ या
$x^{\nssucc\dots\nssucc} = y$। यदि $y^{\nssucc\dots\nssucc} = x$
(जिसमें $n$ बार $\nssucc$ है), तो $y \nsplus n = x$, और इसका विलोम भी,
क्योंकि $\Th{PA} \Proves \lforall[x][\eq[x^{\prime\dots\prime}][(x +
    \num{n})]]$ (यदि $n$,~$\prime$ की संख्या है)।
\end{proof}

\begin{prop}
यदि $[x] \neq [y]$ और $x \nsless y$, तो किसी भी $u \in [x]$ और किसी भी
$v \in [y]$ के लिए $u \nsless v$।
\end{prop}

\begin{proof}
ध्यान दें कि $\Th{PA} \Proves \lforall[x][\lforall[y][(x < y \lif (x' < y
    \lor x' = y))]]$। अतः यदि $u \nsless v$, तो
$u \nsplus
n^\nssucc \nsless v$ भी किसी~$n$ के लिए, यदि $[u] \neq [v]$।

प्रत्येक $u \in [x]$,~$\nsless y$ है: परिकल्पना से $x \nsless y$। यदि $u
\nsless x$, तो संक्रामकता से $u \nsless y$। और यदि $x \nsless u$, किंतु
$u
\in [x]$, तो हमारे पास $u = x\nsplus n^\nssucc$ है किसी~$n$ के लिए, और
इसलिए अभी सिद्ध किए गए तथ्य से $u
\nsless y$।

अब मान लें कि कोई $v \in [y]$,~$\nsless y$ है, अर्थात्
$v \nsplus
m^\nssucc = y$ किसी मानक~$m$ के लिए। इससे $v \nsless x$ असंभव हो जाता है,
क्योंकि अन्यथा $y = v \nsplus m^\nssucc \nsless x$। स्पष्टतः
$x \neq
v$ भी, क्योंकि अन्यथा
$x \nsplus m^\nssucc = v \nsplus m^\nssucc = y$ और हमें
$[x] = [y]$ मिलता। अतः $x \nsless v$। किंतु तब
$x \nsplus
n^\nssucc \nsless v$ भी किसी भी~$n$ के लिए। इसलिए यदि $x \nsless u$ और
$u \in
[x]$, तो $u \nsless v$। यदि $u \nsless x$, तो संक्रामकता से
$u \nsless v$।

अंततः यदि $y \nsless v$, तो $u \nsless v$, क्योंकि जैसा हमने दिखाया है, $u
\nsless y$ और $y \nsless v$।
\end{proof}

\begin{cor}
यदि $[x] \neq [y]$, तो $[x] \cap [y] = \emptyset$।
\end{cor}

\begin{proof}
मान लें कि $z \in [x]$ और $x \nsless y$। तब $z \nsless u$ सभी $u
\in [y]$ के लिए। यदि $z \in [y]$, तो $z \nsless z$ होता। $y
\nsless x$
की स्थिति में भी इसी प्रकार।
\end{proof}

\begin{explain}
इसका अर्थ है कि स्वयं खंडों को ऐसे क्रमित किया जा सकता है जो $\nsless$ का
सम्मान करे: $[x] \nsless [y]$ तभी और केवल तभी जब $x \nsless y$, अथवा
समतुल्य रूप से, जब $u \nsless v$ किसी भी $u \in [x]$ और $v \in [y]$ के
लिए। स्पष्टतः मानक खंड $[0]$ सबसे छोटा खंड है। वह किसी अमानक खंड
को प्रतिच्छेद नहीं करता, और कोई दो अमानक खंड भी परस्पर प्रतिच्छेद नहीं करते।
विशेषतः, बार-बार परवर्ती या पूर्ववर्ती लेकर किसी दूसरे खंड तक “पहुँचा” नहीं
जा सकता।
\end{explain}

\begin{prop}
यदि $x$ और $y$ अमानक हैं, तो $x \nsless x \nsplus y$ और $x
\nsplus y \notin [x]$।
\end{prop}

\begin{proof}
यदि $y$ अमानक है, तो $y \neq \nszero$। $\Th{PA} \Proves
\lforall[x][(y \neq \Obj{0} \lif x < (x+y))]$। अब मान लें कि $x
\nsplus y \in [x]$। चूँकि $x \nsless x \nsplus y$, हमें $x
\nsplus n^\nssucc = x \nsplus y$ मिलता। किंतु
$\Th{PA} \Proves
\lforall[x][\lforall[y][\lforall[z][(\eq[(x+y)][(x+z)] \lif y = z)]]]$
(योग के लिए निरसन नियम)। इसका अर्थ होता कि $y = n^\nssucc$ किसी मानक~$n$
के लिए; किंतु $y$ को अमानक माना गया है।
\end{proof}

\begin{prop}
कोई सबसे छोटा अमानक खंड नहीं है।
\end{prop}

\begin{proof}
$\Th{PA} \Proves \lforall[x][\lexists[y][(\eq[(y+y)][x] \lor
      \eq[(y+y)'][x])]]$; अर्थात् प्रत्येक $x$,~$2$ से विभाज्य है
  (संभवतः शेषफल~$1$ के साथ)। यदि $x$ अमानक है, तो $y$ भी अमानक है। पिछले
  प्रस्ताव से $y \nsless y \nsplus y$ और $y \nsplus y
  \notin [y]$। तब $y \nsless (y \nsplus y)^\nssucc$ और $(y
  \nsplus y)^\nssucc \notin [y]$ भी। किंतु
  $x = y \nsplus y$ या $x = (y
  \nsplus y)^\nssucc$, इसलिए
  $y \nsless x$ और $y \notin [x]$।
\end{proof}

\begin{prop}
कोई सबसे बड़ा खंड नहीं है।
\end{prop}

\begin{proof}
अभ्यास।
\end{proof}

\begin{prob}
दिखाइए कि~$\Th{PA}$ के किसी अमानक प्रतिरूप में कोई सबसे बड़ा खंड नहीं है।
\end{prob}

\begin{prop}
\ollabel{prop:blocks-dense}
खंडों का क्रम सघन है। अर्थात् यदि $x \nsless y$ और $[x] \neq [y]$, तो
उनके बीच कोई ऐसा खंड $[z]$ है जो दोनों से भिन्न है।
\end{prop}

\begin{proof}
मान लें कि $x \nsless y$। पहले की तरह $x \nsplus y$ दो से विभाज्य है
(संभवतः शेषफल सहित): कोई $z \in \Domain{M}$ ऐसा है कि या तो
$x \oplus y = z \oplus z$ या $x \oplus y = (z \oplus
z)^\nssucc$। अवयव $z$,~$x$ और~$y$ का “औसत” है, और $x
\nsless z$ तथा
$z \nsless y$।
\end{proof}

\begin{prob}
\olref[mod][mar][mpa]{prop:blocks-dense} का विस्तृत प्रमाण लिखिए।
कौन-सा !!{sentence} $\Th{PA}$ को !!{derive} करना चाहिए ताकि~$z$ का अस्तित्व
प्रत्याभूत हो? $x \nsless z$ और $z \nsless y$ क्यों, तथा
$[x] \neq [z]$ और $[z]
\neq [y]$ क्यों?
\end{prob}

\begin{explain}
इसलिए अमानक खंड परिमेय संख्याओं की तरह क्रमित हैं: वे अंतबिंदु-रहित
!!a{denumerable} सघन रैखिक क्रम बनाते हैं। यह दिखाया जा सकता है कि ऐसे कोई
भी दो !!{denumerable} क्रम समरूपी होते हैं। इससे निष्कर्ष निकलता है
कि सत्य अंकगणित के किन्हीं दो !!{enumerable} अमानक प्रतिरूपों $\Struct{M}_1$ और
$\Struct{M_2}$ के, केवल $<$ और $=$ वाली भाषा तक न्यूनीकृत रूप समरूपी होते
हैं। वास्तव में, एक समरूपता $h$
इस प्रकार परिभाषित की जा सकती है: $\Struct{M_1}$ और $\Struct{M_2}$ के मानक
भाग मानक प्रतिरूप $\Struct{N}$ के समरूपी हैं और इसलिए परस्पर भी समरूपी हैं।
अमानक भाग बनाने वाले खंड स्वयं परिमेय संख्याओं की तरह क्रमित हैं और इसलिए
समरूपी हैं; खंडों की किसी समरूपता को, प्रत्येक खंड के मनमाने अवयवों को आपस
में मिलाकर और फिर~$x$ के परवर्ती के,~$\Struct{M_1}$ में, प्रतिबिंब को
~$x$ के प्रतिबिंब का,~$\Struct{M_2}$ में, परवर्ती लेकर, खंडों के
\emph{भीतर} की समरूपता तक विस्तृत किया जा सकता है। ध्यान दें कि इससे यह
निष्कर्ष \emph{नहीं} निकलता कि $\mathfrak{M}_1$ और $\mathfrak{M}_2$
अंकगणित की पूरी भाषा में समरूपी हैं (वास्तव में, समरूपता सदैव किसी
!!a{language} के सापेक्ष होती है), क्योंकि $\Domain{M_1}$ और
$\Domain{M_2}$ पर योग तथा गुणन को परिभाषित करने के परस्पर असमरूपी ढंग हैं।
(यह वॉट के प्रसिद्ध प्रमेय से भी निकलता है कि किसी पूर्ण सिद्धांत के गणनीय
प्रतिरूपों की संख्या~2 नहीं हो सकती।)
\end{explain}
\end{document}
